\section{Related Work}

\subsection{Traditional Code Translation Methods}

Early work recognized that programs possess rich syntactic structures and that direct application of sequence-to-sequence models often yields syntactically invalid outputs. To address this, Chen et al. \cite{chen2018tree} pioneered the use of tree-to-tree neural networks, which encode and decode abstract syntax trees (ASTs) to explicitly model program structure. This was further refined by Drissi et al. \cite{drissi2018program}, who introduced a grammar-driven tree-to-tree model that hard-codes target language syntax rules into the decoder, guaranteeing grammatical correctness. These studies established the importance of structural awareness in code translation but were limited by their dependence on parsers, manually defined grammars, and parallel tree-aligned corpora, which hindered generalization and scalability.

To overcome the scarcity of parallel data, The field shifted toward data-driven, unsupervised methods. Roziere et al.\cite{roziere2020unsupervised} introduced TransCoder, which leverages back-translation and denoising auto-encoding on monolingual code corpora to learn cross-lingual code representations without supervision. Concurrently, pre-trained models such as CodeBERT and CodeT5 adapted the pre-training paradigm from natural language processing to code, learning general-purpose representations via masked language modeling. These approaches demonstrated the power of scale and transfer learning, moving the focus from designing task-specific architectures to fine-tuning general-purpose models. However, Malyala et al.\cite{malyala2023ml} found that purely data-driven models often exhibit systematic errors in complex translation scenarios, such as type conversions and loop restructuring. Pan et al.\cite{pan2024lost} conducted a large-scale empirical study showing that LLMs introduce a significant number of bugs during translation, particularly in real-world projects, highlighting the limitations of statistical approaches in understanding deep semantic and syntactic constraints.

\subsection{LLM Based Code Translation Methods}

To mitigate the shortcomings of purely data-driven models, recent research emphasizes augmenting LLMs with programming language-specific knowledge, such as type systems, safety rules, and API mappings. This is especially critical for security-sensitive migrations, such as from C to Rust. Luo et al. \cite{luo2025integrating} proposed the IRENE framework, which combines static analysis to extract ``high-risk patterns'' with retrieval-augmented generation to guide the model toward producing Rust code that complies with ownership rules. Similarly, Hong et al. \cite{hong2025type} focused on type migration, employing a call-graph-based strategy to coordinate interface changes across functions. These hybrid systems effectively address the ``semantic drift'' and safety violations common in naive LLM translations. Extending this neuro-symbolic approach, Ibrahimzada et al. \cite{ibrahimzada2025alphatrans} introduced AlphaTrans, which compositionally combines symbolic knowledge with neural models for repository-level translation and validation. Meanwhile, Nitin et al.\cite{nitin2024spectra} enhances translation by generating multi-modal specifications (e.g., comments, tests) to provide additional context. While these methods significantly improve reliability, they often require substantial engineering effort to construct and maintain domain-specific knowledge bases.

\subsection{Hallucination in General-Purpose LLMs}

Hallucination in large language models (LLMs) refers to the generation of factually incorrect or contextually inconsistent content. Early research focused on text-based hallucinations, with recent work expanding to multimodal settings \cite{rawte2023survey}. Detection methods have evolved from lexical metrics (ROUGE, BLEU) to more robust frameworks incorporating multi-dimensional evaluation \cite{chen2023hallucination} and domain-specific approaches like self-reflection in medical QA \cite{ji2023towards}. Mitigation strategies include knowledge augmentation, prompt engineering, and adversarial defenses \cite{yao2023llm,roychowdhury2024journey}.

Code generation introduces unique challenges due to strict syntactic and functional requirements. Liu et al. \cite{liu2026beyond} identified three primary hallucination types: Requirement Conflict, Knowledge Hallucination, and Code Inconsistency. In repository-level generation, Zhang et al.\cite{zhang2025llm} further categorized these into Task Requirement Conflicts, Factual Knowledge Conflicts, and Project Context Conflicts.

The causes are multifaceted, including training data quality issues, intent comprehension limitations, and inadequate context awareness \cite{zhang2025llm}. Impacts are severe: 95.30\% of hallucinations cause functional errors, with hallucinated code having only 6.26\% pass@1 rate versus 54.73\% for clean code \cite{liu2026beyond}.

Retrieval-augmented generation (RAG) has proven effective, improving pass@1 scores across models (e.g., +4.38\% for ChatGPT) \cite{zhang2025llm}. Other approaches include self-refinement prompts, chain-of-thought reasoning, and specialized code models. 

\subsection{Repo-level Code Translation Methods}

As model capabilities have grown, the research scope has expanded from function-level snippets to more complex, practical scenarios. Repository-level code translation presents unique challenges, including managing inter-file dependencies, build configurations, and project-wide context. Wang et al. \cite{wang2024repotransbench} introduced RepoTransBench, a large-scale, multi-language benchmark for evaluating repository-level translation, revealing that even state-of-the-art LLMs struggle (with success rates around 32.8\%). To tackle this, Zhang et al. \cite{zhang2025skeleton} proposed a skeleton-guided translation framework that first extracts a high-level structural skeleton in the target language before filling in implementation details. Yuan et al.\cite{yuan2025project} and Ou et al.\cite{K3Trans2025} further explored project-level translation by synergistically integrating knowledge graphs and LLMs, or by evolving a triple knowledge-augmented LLM within the repository context.

Parallel to scaling up, improving the human-in-the-loop experience has become a key focus. Liu et al.\cite{liu2024hmcodetrans} developed hmCodeTrans, an interactive translation framework that allows developers to provide real-time feedback. By employing attention caching and suffix splicing, the system achieves millisecond-level responsiveness, reducing human effort by approximately 65\% compared to traditional post-editing workflows. This marks a shift from pursuing full automation toward designing collaborative tools that integrate seamlessly into developer workflows.

Our work is situated at this intersection. We aim to develop a translation framework that leverages the generative strength of LLMs while deeply integrating extensible, domain-aware knowledge and validation mechanisms. We target translation tasks involving complex algorithm logic and API migrations (akin to Type 3/4 in the taxonomy by Jiao et al.\cite{jiao2023evaluation}). Inspired by hybrid systems like IRENE \cite{luo2025integrating} and compositional approaches like AlphaTrans \cite{ibrahimzada2025alphatrans}, we propose a more generalized and extensible translation knowledge orchestration method. This approach is designed to dynamically retrieve and apply relevant constraints and patterns, ensuring high semantic fidelity and project-level usability without sacrificing flexibility.
